\section{Why the decision problem is NP-complete}
\subsection*{Checking a proposed solution: membership in NP}
A certificate is the list of routes. Create a Boolean array with one entry per clinic. Scan the
routes, mark their clinics, reject duplicate or invalid IDs, add demands and edge costs. Finally
check that every clinic was seen, each load is at most $Q$, at most $K$ routes are used, and
the cost is at most $B$.

A valid certificate contains $n$ clinic visits and at most $n+K\le2n$ edges, taking $K\le n$.
With a distance matrix already available, verification takes $O(n)$ arithmetic operations and
$O(n)$ extra memory. Thus CVRP decision is in NP. The formal argument uses finitely encoded
rational or integer edge costs. Building the matrix from coordinates is separate $O(n^2)$
preprocessing; the program uses a tolerance for floating-point cost checks.

\subsection*{Step A: Hamiltonian Cycle to TSP decision}
Start with an undirected graph $H$ with $N$ vertices. Hamiltonian Cycle asks whether there is a
cycle visiting every vertex exactly once [1]. Construct a complete graph on the same vertices.
Give an edge cost 1 if it exists in $H$ and cost 2 otherwise. Set the TSP bound to $B=N$.
Constructing these costs examines $O(N^2)$ pairs.

\textbf{Forward direction.} A Hamiltonian cycle of $H$ uses $N$ original edges. Each costs 1 in
the new graph, giving a TSP tour of cost $N$.

\textbf{Reverse direction.} Every TSP tour has $N$ edges costing at least 1. A total cost at most
$N$ forces every edge to cost 1. Those edges all belong to $H$, giving a Hamiltonian cycle.

The construction is metric: a direct edge costs at most 2, while a detour through a distinct
third vertex costs at least $1+1=2$.

\subsection*{Step B: TSP decision to CVRP decision}
Take a TSP instance with $n+1$ vertices. Choose one as depot 0; the other $n$ become clinics.
Copy the distances, set every demand to 1, set $Q=n$, set $K=1$, and keep bound $B$.
Copying the matrix takes $O(n^2)$ time.

A TSP tour becomes a single CVRP route with load $n$, so it is valid and has the same cost.
Conversely, with one vehicle, a valid CVRP solution is a single depot-to-depot tour visiting all
clinics. It gives a TSP tour of the same cost. Both directions preserve the answer.

\subsection*{Why the chain works}
The chain is Hamiltonian Cycle $\le_p$ TSP decision $\le_p$ CVRP decision. A polynomial-time
construction cannot produce more than polynomially many symbols. Applying another polynomial-time
construction to that output still takes polynomial time. Since Hamiltonian Cycle is NP-complete,
CVRP decision is NP-hard. Together with membership in NP, it is NP-complete.

\takeaway{An exact polynomial-time algorithm for every CVRP instance would also solve the
known-hard problem at the start of the chain. This motivates searching for good feasible plans
within a time budget. NP-hardness does not prove that every particular instance is difficult.}
